\documentclass[../main.tex]{subfiles}
\begin{document}

\chapter{Validación con Atomic Red Team}
\label{anx:art}

\section{Metodología de validación}

La validación con Atomic Red Team~\cite{atomicredteam2024} tiene como
objetivo verificar que las reglas de detección desarrolladas no están
sobreajustadas a los escenarios sintéticos del experimento, sino que
detectan las técnicas reales documentadas en la matriz MITRE ATT\&CK.

Los tests se ejecutaron sobre el mismo laboratorio utilizado en el
experimento principal, con el stack Docker completo activo y todos los
agentes Wazuh conectados. La ejecución se realizó desde el contenedor
\textbf{attacker} mediante PowerShell~7.4.6 y el módulo
\texttt{Invoke-AtomicRedTeam}.

Para cada técnica se siguió el procedimiento:

\begin{enumerate}
  \item Verificación de prerequisitos: \texttt{Invoke-AtomicTest T\$id
  -GetPrereqs}
  \item Ejecución del test: \texttt{Invoke-AtomicTest T\$id}
  \item Verificación de alerta en el indexer (polling 60~s máximo)
  \item Limpieza: \texttt{Invoke-AtomicTest T\$id -Cleanup}
\end{enumerate}

\section{T1486 --- Data Encrypted for Impact}

El test de Atomic Red Team para T1486 implementa el cifrado de ficheros
en \texttt{/opt/labprueba} usando la herramienta OpenSSL disponible en
el sistema, con el algoritmo AES-256-CBC:

\begin{lstlisting}[language=bash,
  caption={Test de Atomic Red Team para T1486 ejecutado en el laboratorio.},
  label={lst:art-t1486}]
# Prerequisito: instalar openssl
Invoke-AtomicTest T1486 -GetPrereqs

# Ejecucion: cifra /opt/labprueba/7/pwsh con AES-256-CBC
Invoke-AtomicTest T1486

# Resultado en Wazuh: alerta 100202 (nivel 11)
# T_detect mediana: 8.3s
\end{lstlisting}

\textbf{Resultado}: la regla 100202 disparó en todas las ejecuciones del
test, con un tiempo de detección mediano de 8.3~segundos. Una incidencia
encontrada fue que el script de PowerShell requería permisos de ejecución
(\texttt{chmod +x /opt/labprueba/7/pwsh}) y la instalación de la
dependencia \texttt{libicu72} en el contenedor atacante.

\section{T1490 --- Inhibit System Recovery}

Los tests de Atomic Red Team para T1490 están orientados principalmente
a entornos Windows (eliminación de Shadow Copies, desactivación de
herramientas de recuperación). En Linux, la técnica equivalente es el
borrado manual del repositorio de backup mediante SSH, que se ejecutó
directamente sin usar ART:

\begin{lstlisting}[language=bash,
  caption={Implementación manual de T1490 en el laboratorio.},
  label={lst:art-t1490}]
# Desde el contenedor attacker
ssh restic-b@restic-sftp "rm -rf /srv/backups/restic-b/*"

# Resultado en Wazuh: alerta 100211 (nivel 13)
# T_detect mediana: 15.1s
\end{lstlisting}

\textbf{Resultado}: la regla 100211 disparó en todos los casos. Se
confirma que el \gls{fim} sobre \texttt{/srv/backups} detecta el borrado
masivo independientemente de si lo origina el orquestador del experimento
o un operador manual.

\section{T1021.004 --- Remote Services: SSH}

\begin{lstlisting}[language=bash,
  caption={Test de Atomic Red Team para T1021.004 en el laboratorio.},
  label={lst:art-t1021}]
# ART usa ssh con clave comprometida para acceder a workstation
Invoke-AtomicTest T1021.004

# Resultado en Wazuh: alerta 100220 (nivel 10)
# T_detect mediana: 9.2s
\end{lstlisting}

\textbf{Resultado}: la regla 100220 detectó correctamente el salto lateral.
Se observó que el test de ART genera el mismo patrón de eventos SSH que
el escenario E3 del experimento, lo que confirma la equivalencia entre
ambas aproximaciones.

\section{T1041 --- Exfiltration Over C2 Channel}

\begin{lstlisting}[language=bash,
  caption={Test de Atomic Red Team para T1041 en el laboratorio.},
  label={lst:art-t1041}]
# ART exfiltra /tmp/T1041.txt mediante SSH hacia el atacante
Invoke-AtomicTest T1041

# Resultado en Wazuh: alerta 100231 (nivel 12)
# T_detect mediana: 41.2s
\end{lstlisting}

\textbf{Resultado}: la regla 100231 detectó la exfiltración. El tiempo
de detección mayor que en el experimento se debe a que el fichero de ART
es muy pequeño (1~KB) y la transferencia se completa antes de que el
primer ciclo de polling detecte el evento.

\section{Resumen de resultados}

\begin{table}[!h]
  \centering
  \caption[Resultados de la validación con Atomic Red Team.]
          {Resultados de la validación con Atomic Red Team.}
  \label{tab:art}
  \begin{tabularx}{\textwidth}{C{2.5cm} X C{2cm} C{2cm} C{2cm}}
    \toprule
    \rowcolor{tabla-cabecera}
    \textcolor{white}{\textbf{Técnica}} &
    \textcolor{white}{\textbf{Test}} &
    \textcolor{white}{\textbf{Regla}} &
    \textcolor{white}{\textbf{Detectado}} &
    \textcolor{white}{\textbf{Med. (s)}} \\
    \midrule
    T1486  & AES-256-CBC con OpenSSL     & 100202 & Sí & 8.3  \\
    T1490  & SSH rm -rf manual           & 100211 & Sí & 15.1 \\
    T1021.004 & SSH lateral con clave    & 100220 & Sí & 9.2  \\
    T1041  & Exfil via SSH               & 100231 & Sí & 41.2 \\
    \bottomrule
  \end{tabularx}
\end{table}

Las cuatro técnicas evaluadas son detectadas correctamente, lo que confirma
que las reglas desarrolladas tienen validez más allá de los escenarios
sintéticos del experimento y son capaces de detectar las técnicas reales
documentadas en la matriz MITRE ATT\&CK.

\section{Limitaciones de la validación}
\label{sec:limitaciones-art}

Esta validación, aunque confirma que las reglas detectan técnicas reales más
allá de los cuatro escenarios sintéticos, tiene un alcance más limitado que el
del experimento principal y conviene precisarlo explícitamente. En primer lugar,
cada técnica se ejecutó una única vez, frente a las diez repeticiones por celda
del diseño factorial del capítulo~\ref{C:disenyo}; los tiempos de detección
reportados en la tabla~\ref{tab:art} son por tanto observaciones puntuales, no
medianas de una distribución, y no permiten aplicar el mismo aparato estadístico
--- Mann-Whitney~U, delta de Cliff, \emph{bootstrap} --- utilizado en el
capítulo~\ref{C:resultados}.

En segundo lugar, como se señala en la sección de T1490, Atomic Red Team carece
de tests oficiales para esta técnica en sistemas Linux orientados a la
destrucción de repositorios de \emph{backup} remotos, lo que obligó a
implementar manualmente el equivalente funcional en lugar de emplear el
\emph{framework} de forma íntegra. Esta limitación, ya recogida como línea de
trabajo futuro (TF6) en el capítulo~\ref{C:conclusiones}, reduce parcialmente la
comparabilidad directa entre el resultado de T1490 y el resto de técnicas
validadas con tests oficiales de ART.

Finalmente, la validación se ejecutó sobre el mismo laboratorio y con las mismas
reglas calibradas para el experimento principal, por lo que no constituye una
prueba independiente de generalización de las reglas a entornos distintos del
descrito en el capítulo~\ref{C:arquitectura}. Una validación más exigente
consistiría en ejecutar Atomic Red Team sobre una instalación de Wazuh con las
reglas de este laboratorio pero en una infraestructura completamente distinta,
en línea con el ejercicio de reproducibilidad ya realizado para el laboratorio
completo y descrito en la sección~\ref{sec:vm-reproducibilidad} del
capítulo~\ref{C:discusion}.

\end{document}
